% UChicago-style poster draft for MACS 30200
% Three-column version designed to fill the 48 x 36 poster page.
\documentclass[final]{beamer}

\usepackage[T1]{fontenc}
\usepackage{lmodern}
\usepackage[size=custom,width=121.92,height=91.44,scale=1.12]{beamerposter}
\usetheme{gemini}
\usecolortheme{uchicago}
\usepackage{graphicx}
\usepackage{booktabs}
\usepackage{tabularx}
\usepackage[patch=none]{microtype}
\usepackage{tikz}
\usepackage{amsmath}
\usepackage{hyperref}

\title{Does Remote Work Reduce or Reproduce Inequality in the AI Labor Market?}
\author{Qinghao Zhong}
\institute{MACS 30200 Computational Research Design \quad University of Chicago}

\newlength{\sepwidth}
\newlength{\colwidth}
\setlength{\sepwidth}{0.025\paperwidth}
\setlength{\colwidth}{0.30\paperwidth}
\newcommand{\separatorcolumn}{\begin{column}{\sepwidth}\end{column}}

\footercontent{
  \href{https://github.com/qinghaoz-cell/macs30200}{github.com/qinghaoz-cell/macs30200} \hfill
  Data: ai-jobs.net / Kaggle Salaries for Data Science Jobs \hfill
  Qinghao Zhong
}

\begin{document}

\begin{frame}[t]
\begin{columns}[t]
\separatorcolumn

% ============================================================
% Column 1: question, motivation, data, methods
% ============================================================
\begin{column}{\colwidth}

  \begin{alertblock}{Research Question}
    Does remote work reduce or reproduce salary inequality in AI- and data-related jobs?

    \smallskip
    This project asks whether remote work is associated with more equal salary outcomes, or whether it remains tied to existing labor-market sorting.
  \end{alertblock}

  \begin{block}{Background \& Motivation}
    Remote work is often described as flexible, global, and opportunity-expanding. But research on remote work shows that access to remote jobs is uneven across occupations, workers, firms, and places.

    \smallskip
    AI and data jobs are a useful case because they are high-skill, digitally mediated, and commonly imagined as location-flexible. If remote work is an equalizer, it should be visible in this type of labor market.

    \smallskip
    \textbf{Contribution:} I connect remote-work debates to salary inequality by testing whether flexibility translates into more equal pay.
  \end{block}

  \begin{block}{Data Summary}
    \begin{itemize}
      \item Public salary data for AI, machine-learning, and data-related jobs.
      \item Source: Kaggle dataset based on ai-jobs.net salary submissions.
      \item Analytic sample: full-time jobs only, \textbf{150,541 observations}.
      \item Key variables: salary in USD, work category, experience level, company location, company size, job title, and work year.
    \end{itemize}

    \begin{center}
    \renewcommand{\arraystretch}{1.35}
    \LARGE
    \begin{tabular}{lrr}
      \toprule
      \textbf{Category} & \textbf{N} & \textbf{Mean} \\
      \midrule
      On-site & 118,972 & \$159,716 \\
      Remote & 31,256 & \$151,887 \\
      Hybrid & 313 & \$82,814 \\
      \bottomrule
    \end{tabular}
    \end{center}
  \end{block}

  \begin{block}{Methods and Measures}
    \begin{itemize}
      \item Recoded remote ratio into on-site, hybrid, and remote categories.
      \item Compared raw salary distributions across work categories.
      \item Compared salary patterns across top company locations.
      \item Matched remote and on-site jobs within the same experience level, company location, company size, and job title.
      \item Estimated adjusted salary associations using log salary as the outcome.
    \end{itemize}

    \begin{center}
    \begin{minipage}{0.92\linewidth}
    \renewcommand{\arraystretch}{1.22}
    \Large
    \begin{tabularx}{\linewidth}{lX}
      \toprule
      \textbf{Concept} & \textbf{Operational measure} \\
      \midrule
      Salary outcome & Annual salary in USD \\
      Work category & Remote ratio: 0, 50, or 100 \\
      Firm context & Company size and location \\
      Job similarity & Job title and experience match \\
      \bottomrule
    \end{tabularx}
    \end{minipage}
    \end{center}
  \end{block}

  \begin{block}{Working Expectations}
    \begin{description}
      \item[\textbf{H1}] If remote work equalizes opportunity, remote jobs should have similar or higher salaries than on-site jobs.
      \item[\textbf{H2}] If remote work reproduces sorting, salary differences should remain after accounting for job and firm characteristics.
    \end{description}
  \end{block}

\end{column}
\separatorcolumn

% ============================================================
% Column 2: main findings and descriptive results
% ============================================================
\begin{column}{\colwidth}

  \begin{alertblock}{Main Finding}
    The evidence does \textbf{not} support a simple story that remote work equalizes salary outcomes. Remote work appears tied to broader job, location, firm, and experience-based sorting.
  \end{alertblock}

  \begin{block}{Result 1: Salary Distributions}
    \begin{figure}
      \centering
      \includegraphics[width=\linewidth,height=0.265\paperheight,keepaspectratio]{salary_distribution_with_means.png}
    \end{figure}

    \textbf{Finding:} Remote and on-site jobs both occupy the high-salary portion of the market, but remote jobs do not clearly exceed on-site jobs in the raw comparison. Hybrid jobs are a small and lower-paid group in this dataset.
  \end{block}

  \begin{block}{Result 2: Geography Matters}
    \begin{figure}
      \centering
      \includegraphics[width=\linewidth,height=0.255\paperheight,keepaspectratio]{Figures/Figure_3.png}
    \end{figure}

    \textbf{Finding:} Salary differences across work categories vary across top company locations. This suggests that remote work is embedded in local and firm-specific labor markets rather than operating uniformly everywhere.
  \end{block}

  \begin{block}{Interpretive Takeaway}
    \begin{itemize}
      \item Flexibility is not the same as equality.
      \item Remote work may expand where people can work, but salary returns still depend on job type, geography, firm context, and experience.
    \end{itemize}
  \end{block}

\end{column}
\separatorcolumn

% ============================================================
% Column 3: matched results, implications, future work
% ============================================================
\begin{column}{\colwidth}

  \begin{block}{Result 3: Matched Comparison}
    \begin{figure}
      \centering
      \includegraphics[width=\linewidth,height=0.27\paperheight,keepaspectratio]{matched_remote_onsite_salary_clean.png}
    \end{figure}

    \textbf{Finding:} Comparing remote and on-site observations within the same experience level, company location, company size, and job title, remote jobs remain lower on average in the matched sample.
  \end{block}

  \begin{block}{Result 4: Adjusted Association}
    \begin{center}
    \renewcommand{\arraystretch}{1.25}
    \Large
    \begin{tabularx}{\linewidth}{lcl}
      \toprule
      \textbf{Comparison} & \textbf{Est.} & \textbf{Meaning} \\
      \midrule
      Remote vs. on-site & -5.9\% & lower salary \\
      Hybrid vs. on-site & -12.0\% & lower salary \\
      \bottomrule
    \end{tabularx}
    \end{center}

    Model controls include experience level, company size, work year, and company location. These estimates describe associations, not causal effects.
  \end{block}

  \begin{block}{Contributions \& Implications}
    \begin{enumerate}
      \item \textbf{Flexibility is not wage equality.} Remote AI/data jobs are not automatically associated with higher pay.
      \item \textbf{Remote work is structured by sorting.} Salary returns still vary with job type, firm context, geography, and experience.
      \item \textbf{Employer implication.} Claims about remote work should separate access, flexibility, and wage equality.
    \end{enumerate}
  \end{block}

  \begin{block}{Limitations and Future Research}
    \begin{itemize}
      \item Observational analysis; results should not be interpreted causally.
      \item Dataset lacks worker demographics, household constraints, preferences, and detailed skills.
      \item Hybrid jobs are rare, so hybrid estimates are less stable than remote/on-site comparisons.
      \item Next steps: stronger job-title groupings, alternative matching specifications, and additional robustness checks.
    \end{itemize}
  \end{block}

  \begin{block}{Code and Data}
    \footnotesize
    \href{https://github.com/qinghaoz-cell/macs30200}{github.com/qinghaoz-cell/macs30200}
  \end{block}

  \begin{block}{References}
    \scriptsize
    \begin{itemize}
      \item Aksoy, C. G., Barrero, J. M., Bloom, N., Davis, S. J., Dolls, M., \& Zarate, P. (2022). Working from home around the world. \textit{Brookings Papers on Economic Activity}, \textit{Fall 2022}, 281--360.
      \item Cetrulo, A., Guarascio, D., \& Virgillito, M. E. (2022). Working from home and the explosion of enduring divides: Income, employment and safety risks. \textit{Structural Change and Economic Dynamics}, \textit{60}, 244--258.
      \item Dingel, J. I., \& Neiman, B. (2020). How many jobs can be done at home? \textit{Journal of Public Economics}, \textit{189}, 104235.
      \item Hansen, S., Lambert, P. J., Bloom, N., Davis, S. J., Sadun, R., \& Taska, B. (2023). \textit{Remote work across jobs, companies, and space}. IZA Discussion Paper No. 16420.
      \item Lee, K. (2023). Working from home as an economic and social change: A review. \textit{Labour Economics}, \textit{85}, 102462.
      \item Pabilonia, S. W., \& Vernon, V. (2024). Remote work, wages, and hours worked in the United States. \textit{Monthly Labor Review}.
    \end{itemize}
  \end{block}

\end{column}
\separatorcolumn
\end{columns}
\end{frame}

\end{document}
